%! Author = denis
%! Date = 11/08/2026

\begin{frame}{Identidades lógicas: Lei da Absorção}
    \small

    \begin{block}{Definição}
        Uma variável pode ``absorver'' uma expressão que já contém 
        essa mesma variável.
    \end{block}

    \begin{columns}[T]

        % Primeira forma
        \small
        \begin{column}{0.49\textwidth}
            \begin{block}{Absorção com OR}
                \[
                    \boxed{A + AB = A}
                \]

                Demonstração:
                \[
                    A + AB
                \]
                \[
                    = A(1+B)
                \]
                \[
                    = A\cdot1 = \boxed{A}
                \]
            \end{block}
        \end{column}

        % Segunda forma
        \small
        \begin{column}{0.49\textwidth}
            \begin{block}{Absorção com AND}
                \[
                    \boxed{A(A+B)=A}
                \]

                Demonstração:
                \[
                    A(A+B)
                \]
                \[
                    = AA+AB
                \]
                \[
                    = A+AB = \boxed{A}
                \]
            \end{block}
        \end{column}

    \end{columns}

\end{frame}